\scanpage{47}
% archive sheet number: 46
% The first draft, through ``i.e. un'', is cancelled by long diagonal strokes.
\deleted{Ceci dit, considérons la $r=(j-i-2)$-carte \inserted{$\Pi_{i,j}$} définie par la moquette $(0\leq i\leq i+2\leq j\leq n)$}
\[
\deleted{R_{ij}=R/(\sigma_0\cdots\sigma_{i-1},\sigma_{j+1},\ldots,\sigma_n)}
\]
\deleted{(\uline{NB}}
\[
\deleted{R_{0j}=R/(\sigma_{j+1},\ldots,\sigma_n),\qquad R_{i,n}=R/(\sigma_0\cdots\sigma_{i-1}))}
\]
\deleted{avec les opérations induites par passage au quotient par}
\diagram{$\begin{array}{ccc}
\deleted{\sigma_{i+1}}&\deleted{,\ldots,}&\deleted{\sigma_{j-1}}\\
\deleted{\Vert}&&\deleted{\Vert}\\
\deleted{\sigma'_0}&\deleted{,\ldots,}&\deleted{\sigma'_{j-i-2=r}}
\end{array}$}
\deleted{qui en font une $r$-carte, dont l'ens des composantes connexes s'identifie à $D_j(\Pi)$. Fixons une telle composante connexe, i.e. un}

Ceci dit, soit $f_i\in D_i(\Pi)$ \inserted{$(1\leq i\leq n)$}, $E_i=\unsure{\varphi|\phi}_i(f_i)$, et considérons la $(i-1)$-carte induite $\unsure{\partial|\delta}f_i^{\unsure{\subset|c}}$ dont l'ens des repères est $D_{01\ldots i}(f_i)$ = image rec. de $\{f_i\}$ par
\diagram{\begin{tikzcd}[ampersand replacement=\&]
D_{0,1,\ldots,i}\inserted{(\Pi)} \arrow[r,phantom,"{\unsure{-|\longleftarrow|\longrightarrow}}" description] \& D_i(\Pi).
\end{tikzcd}}
On trouve alors \inserted{en termes de $\varphi$} une réalisation géom. \unsure{canoniq.} (dite ``induite'') de $\partial f_i$ dans $E_i$.

Légèrement moins évidente est la définition de la représentation induite de la $(n-i-1)$-carte $\unsure{\partial f_i^{\supset}|\partial f_i\supseteq}$ dont l'ens des
